\documentclass[../../book]{subfiles}

\begin{document}

%--------------------------------
\subsection{The FRI Protocol}\label{subsec:fri}
%--------------------------------

In \Cref{fig:poly-iop-third-attempt}, the whole compilation of a Poly-IOP into an IOP was based on a single black box: a $\delta$-RS-IOPP (\Cref{def:rs-proximity}) that tests whether one function $h: \Omega \to \mathbb{F}$ is close to $\mathsf{RS}[\mathbb{F}, \Omega, k]$.
In this subsection, we open this ``box''.

Recall why the problem is not trivial.
Without the help of the prover, the verifier needs $k + 1$ queries: it interpolates $h$ through $k$ points and checks the result at one more point.
In applications, $k$ is the size of the whole computation, so such a verifier is as slow as recomputing everything.
The \emph{Fast Reed-Solomon IOP of Proximity} (FRI)~\cite{bensasson18} lets the prover send a few additional oracles, after which the verifier needs only $O(\log n)$ queries.

In fact, we already have the main ingredient.
By \Cref{lemma:two-way-folding}, folding maps the problem ``is $g$ close to $\mathsf{RS}[\mathbb{F}, \Omega, k]$?'' to the same problem for $\mathsf{Fold}(g, r)$, a code of half the dimension, and a domain of half the size.
FRI applies it again and again: after $\log_2 k$ folds the dimension drops to $1$, the code consists of constant functions, and the prover simply sends the constant.
Throughout this subsection, we keep the setting of \Cref{lemma:two-way-folding} and assume that $k = 2^\ell$.
To avoid confusion with the evaluation points $r_i$ of \Cref{fig:poly-iop-third-attempt}, we denote the folding challenges by $\beta$.

%--------------------------------
\subsubsection{Folding Is Local}\label{subsubsec:fri-folding}
%--------------------------------

There is one obstacle: the verifier cannot compute $\mathsf{Fold}(g, \beta)$ itself, since that would require reading all of $g$.
Instead, the prover will send the fold as a new oracle, and the verifier will spot-check it.
This works because folding is \emph{local}: by \Cref{def:fold}, the value of $\mathsf{Fold}(g, \beta)$ at $x^2$ depends only on the two values $g(x)$ and $g(-x)$.
So to check one value of the fold, the verifier reads just two cells of $g$.
There is also a pleasant geometric way to see this, which is how folding is described in~\cite{bensasson18} and illustrated in \Cref{fig:fri-fold-line}.

\begin{exercise}\label{ex:fri-line}
    Let $L$ be the line through the points $(x, g(x))$ and $(-x, g(-x))$, that is, the unique polynomial of degree at most $1$ with $L(x) = g(x)$ and $L(-x) = g(-x)$.
    Show that $\mathsf{Fold}(g, \beta)(x^2) = L(\beta)$.
\end{exercise}

\begin{figure}[ht]
    \centering
    \begin{tikzpicture}[
        lab/.style={font=\footnotesize},
        pt/.style={circle, inner sep=1.7pt},
    ]
        \draw[-{Stealth[length=2mm]}, oc-gray-6] (-2.6, 0) -- (3.6, 0) node[lab, right] {$T$};
        \draw[-{Stealth[length=2mm]}, oc-gray-6] (0, -0.3) -- (0, 2.9);
        \draw[very thick, oc-blue-6] (-2.5, 0.35) -- (3.3, 2.73);
        \node[pt, fill=oc-blue-8] (a) at (-1.5, 0.76) {};
        \node[pt, fill=oc-blue-8] (b) at (1.5, 1.99) {};
        \node[pt, fill=oc-orange-7] (c) at (2.7, 2.48) {};
        \draw[densely dashed, oc-gray-6] (-1.5, 0) node[lab, below] {$-x$} -- (a);
        \draw[densely dashed, oc-gray-6] (1.5, 0) node[lab, below] {$x$} -- (b);
        \draw[densely dashed, oc-orange-7] (2.7, 0) node[lab, below, text=oc-orange-9] {$\beta$} -- (c);
        \node[lab, above left] at (a) {$g(-x)$};
        \node[lab, above left] at (b) {$g(x)$};
        \node[lab, right=2pt, text=oc-orange-9] at (c) {$\mathsf{Fold}(g, \beta)(x^2)$};
        \node[lab, text=oc-blue-8, rotate=22] at (-0.4, 1.45) {line $L$};
    \end{tikzpicture}
    \caption{One value of the fold. The two values of $g$ at the pair $\pm x$ determine a line $L$, and the fold at $x^2$ is the value of this line at the challenge $\beta$ (\Cref{ex:fri-line}). The picture is drawn over the reals only for intuition.}
    \label{fig:fri-fold-line}
\end{figure}

Let us consider some concrete example.

\begin{example}\label{ex:fri-fold}
    Let $\mathbb{F} = \mathbb{F}_{17}$, $\Omega = \langle 2 \rangle$ and $k = 4$, so that $\rho = 1/2$.
    Listing $\Omega$ as $1, 2, 4, 8, 16, 15, 13, 9$ (the powers of $\omega = 2$), the polynomial $f = 3 + T + 2T^2 + 5T^3$ gives the codeword
    \begin{equation*}
        g = \overline{f} = (11, 2, 2, 13, 16, 3, 0, 11).
    \end{equation*}
    Take $\beta = 2$, and recall the circle of \Cref{fig:fri-domains}.
    The first entry of the fold, at $1^2 = 1$, uses the pair $\{1, 16\} = \{1, -1\}$:
    \begin{equation*}
        \mathsf{Fold}(g, 2)(1) = \frac{11 + 16}{2} + 2 \cdot \frac{11 - 16}{2} = 27 \cdot 9 - 5 = 0 \pmod{17},
    \end{equation*}
    where we used $2^{-1} = 9$ in $\mathbb{F}_{17}$.
    Doing the same for the other three pairs gives $\mathsf{Fold}(g, 2) = (0, 2, 10, 8)$ on $\Omega^2$, listed as $1, 4, 16, 13$.
    As expected, these are the evaluations of $f_E + 2 f_O = 5 + 12T$, a polynomial of degree less than $2$.
\end{example}

%--------------------------------
\subsubsection{The Protocol}\label{subsubsec:fri-protocol}
%--------------------------------

To fold $\ell$ times, write $\Omega_i := \Omega^{2^i}$ for the domain after $i$ folds (so $\Omega_0 = \Omega$ and $\Omega_1 = \Omega^2$), and $\code_i := \mathsf{RS}[\mathbb{F}, \Omega_i, k/2^i]$ for the corresponding code; all these codes have the same rate $\rho$, and the last one, $\code_\ell$, consists of the constant functions.
The prover sends every fold as a new oracle, and the verifier spot-checks consecutive oracles.
As in~\cite{bensasson18}, this happens in two phases: in the \emph{commit phase}, the prover sends all oracles, and only afterwards, in the \emph{query phase}, the verifier makes its queries.

\begin{definition}[FRI]\label{def:fri}
    The input is an oracle $g_0 := g: \Omega_0 \to \mathbb{F}$, and the goal is to test proximity to $\code_0 = \mathsf{RS}[\mathbb{F}, \Omega, k]$.
    \begin{itemize}
        \item \textbf{Commit phase.} For every $i \in [\ell]$:
        \begin{algoen}
            \item the verifier samples $\beta_i \leftarrowS \mathbb{F}$;
            \item the prover sends an oracle $g_{i+1}: \Omega_{i+1} \to \mathbb{F}$; an honest prover sets $g_{i+1} := \mathsf{Fold}(g_i, \beta_i)$.
        \end{algoen}
        Finally, the prover sends a constant $c \in \mathbb{F}$, and we set $g_\ell(y) := c$ for all $y \in \Omega_\ell$.
        \item \textbf{Query phase.} Repeat $s$ times:
        \begin{algoen}
            \item the verifier samples $x_0 \leftarrowS \Omega_0$ and sets $x_{i+1} := x_i^2$ for every $i \in [\ell]$;
            \item for every $i \in [\ell]$, the verifier queries $g_i(x_i)$ and $g_i(-x_i)$ and checks the \emph{round consistency}
            \begin{equation*}
                g_{i+1}(x_{i+1}) = \frac{g_i(x_i) + g_i(-x_i)}{2} + \beta_i \cdot \frac{g_i(x_i) - g_i(-x_i)}{2x_i}.
            \end{equation*}
        \end{algoen}
    \end{itemize}
    The verifier accepts if and only if all checks pass.
\end{definition}

Note that the left-hand side of the round consistency check needs no extra query: $g_{i+1}(x_{i+1})$ is one of the two values queried in the next round, and in the last round it is the constant $c$.
The whole protocol is depicted in \Cref{fig:fri-protocol}, and \Cref{fig:fri-tree} below shows the same run as a tree.
By the first claim of \Cref{lemma:two-way-folding}, an honest prover is always accepted: every $g_i$ is a codeword of $\code_i$, so $g_\ell$ is indeed constant.

\begin{figure}[ht]
    \centering
    \scalebox{0.8}{%
    \begin{tikzpicture}[
        cell/.style={draw=oc-gray-7, minimum size=0.32cm, inner sep=0pt},
        qcell/.style={draw=oc-orange-8, thick, fill=oc-orange-3, minimum size=0.32cm, inner sep=0pt},
        send/.style={-{Stealth[length=3mm]}, line width=0.4mm},
    ]
        \node[inner sep=0pt, align=center] (prover) {\includegraphics[width=1.25cm]{contents/images/common/prover.png}\\Prover $\mathcal{P}$};
        \node[inner sep=0pt, align=center, right=7cm of prover] (verifier) {\includegraphics[width=1.25cm]{contents/images/common/verifier.png}\\Verifier $\mathcal{V}^{\textcolor{oc-red-8}{g_0}}$};
        \coordinate (ps) at (prover.south);
        \coordinate (vs) at (prover.south -| verifier.south);
        \path (ps) -- (vs) coordinate[midway] (c);

        \draw [dashed,line width=0.3mm] ([yshift=-0.5cm]ps) -- ([yshift=-9.3cm]ps);
        \draw [dashed,line width=0.3mm] ([yshift=-0.5cm]vs) -- ([yshift=-9.3cm]vs);

        % a strip of #2 cells centred at #1, named #3, fill #4, queried cells #5
        \newcommand{\strip}[5]{
            \foreach \i in {1,...,#2} {
                \node[cell, fill=#4] at ([xshift=0.32cm*\i-0.16cm*#2-0.16cm]#1) {};
            }
            \foreach \i in {#5} {
                \node[qcell] at ([xshift=0.32cm*\i-0.16cm*#2+0.16cm]#1) {};
            }
            \node[left] at ([xshift=-0.16cm*#2-0.1cm]#1) {#3 $=$};
        }
        \newcommand{\fromverifier}[2]{
            \draw[send] ([yshift=#1]vs) -- ([yshift=#1]ps) node[midway, above=0mm]{#2};
        }

        % input oracle
        \node[draw=oc-red-6, dashed, thick, rounded corners=2pt, minimum width=6.6cm, minimum height=1.15cm] at ([yshift=-1.05cm]c) {};
        \node[text=oc-red-8, font=\small] at ([yshift=-0.7cm]c) {input oracle};
        \coordinate (s) at ([yshift=-1.25cm]c);
        \strip{s}{16}{$\textcolor{oc-red-8}{g_0}$}{oc-red-1}{3,11}

        % commit phase
        \fromverifier{-2.5cm}{Sample $\beta_0 \leftarrowS \mathbb{F}$}
        \coordinate (s) at ([yshift=-3.25cm]c);
        \strip{s}{8}{$g_1$}{oc-blue-1}{3,7}
        \draw[send] ([yshift=-3.6cm]ps) -- ([yshift=-3.6cm]vs);
        \fromverifier{-4.7cm}{Sample $\beta_1 \leftarrowS \mathbb{F}$}
        \coordinate (s) at ([yshift=-5.45cm]c);
        \strip{s}{4}{$g_2$}{oc-blue-1}{1,3}
        \draw[send] ([yshift=-5.8cm]ps) -- ([yshift=-5.8cm]vs);
        \node at ([yshift=-6.2cm]c) {$\vdots$};
        \fromverifier{-7.3cm}{Sample $\beta_{\ell-1} \leftarrowS \mathbb{F}$}
        \draw[send] ([yshift=-8.4cm]ps) -- ([yshift=-8.4cm]vs) node[midway, above=0mm]{Send constant $c$};

        % query phase
        \node[align=center, fill=white, below=8.7cm of verifier]{
        \noindent\rule{4.6cm}{0.8pt}\\
        Repeat $s$ times:\\
        sample $x_0 \leftarrowS \Omega_0$, set $x_{i+1} := x_i^2$,\\
        query $g_i(\pm x_i)$, check round consistency\\
        \noindent\rule{4.6cm}{0.8pt}};
    \end{tikzpicture}%
    }
    \caption{The FRI protocol. In the commit phase, the prover answers every challenge $\beta_i$ with the next oracle $g_{i+1}$, which is half as long as $g_i$, and finally sends a constant. In the query phase, the verifier follows a random point down the layers: if $x_0$ is the $4$th cell of $g_0$, then $x_1 = x_0^2$ and $x_2 = x_0^4$ are the $4$th cells of $g_1$ and $g_2$, and the orange cells are the pairs $\pm x_i$ it reads.}
    \label{fig:fri-protocol}
\end{figure}

\begin{example}\label{ex:fri-run}
    We continue \Cref{ex:fri-fold}, where $k = 4$, so FRI has $\ell = 2$ rounds.
    After $\beta_0 = 2$, the honest prover sends $g_1 = (0, 2, 10, 8)$ on $\Omega_1 = \{1, 4, 16, 13\}$.
    After $\beta_1 = 5$, it folds $g_1$ once more and obtains $g_2 = (14, 14)$ on $\Omega_2 = \{1, 16\}$, which is indeed constant, so it sends $c = 14$.
    If the verifier samples $x_0 = 2$, then $x_1 = 4$ and it checks
    \begin{equation*}
        g_1(4) = \frac{g_0(2) + g_0(15)}{2} + 2 \cdot \frac{g_0(2) - g_0(15)}{2 \cdot 2} = \frac{2 + 3}{2} + \frac{2 - 3}{2} = 2,
    \end{equation*}
    and similarly $c = g_1(4)/2 + g_1(13)/2 + 5 \cdot (g_1(4) - g_1(13))/(2 \cdot 4) = 14$.
    The whole run is shown in \Cref{fig:fri-tree}.
\end{example}

\begin{figure}[ht]
    \centering
    \begin{tikzpicture}[
        cell/.style={draw=oc-green-8, fill=oc-green-1, minimum width=0.62cm, minimum height=0.48cm, inner sep=0pt, font=\small},
        hot/.style={draw=oc-orange-8, very thick, fill=oc-orange-2},
        pt/.style={font=\scriptsize, text=oc-gray-6},
        lab/.style={font=\footnotesize},
        edge/.style={oc-green-5, thick},
        hedge/.style={oc-orange-7, very thick},
    ]
        \def\dx{0.98}
        % layer 0: pairs {x, -x} are neighbours
        \foreach \p/\v [count=\j from 0] in {1/11, 16/16, 4/2, 13/0, 2/2, 15/3, 8/13, 9/11} {
            \node[cell] (a\j) at (\dx*\j, 0) {$\v$};
            \node[pt, above=0pt] at (a\j.north) {$\p$};
        }
        % layer 1: the j-th cell sits below the j-th pair
        \foreach \p/\v [count=\j from 0] in {1/0, 16/10, 4/2, 13/8} {
            \node[cell] (b\j) at (\dx*2*\j+\dx/2, -1.7) {$\v$};
            \node[pt, above=0pt] at (b\j.north) {$\p$};
        }
        % layer 2
        \foreach \p/\v [count=\j from 0] in {1/14, 16/14} {
            \node[cell] (c\j) at (\dx*4*\j+\dx*1.5, -3.4) {$\v$};
            \node[pt, above=0pt] at (c\j.north) {$\p$};
        }
        % folding edges
        \foreach \j in {0,...,3} {
            \pgfmathtruncatemacro{\l}{2*\j}\pgfmathtruncatemacro{\r}{2*\j+1}
            \draw[edge] (a\l.south) -- (b\j.north west);
            \draw[edge] (a\r.south) -- (b\j.north east);
        }
        \foreach \j in {0,1} {
            \pgfmathtruncatemacro{\l}{2*\j}\pgfmathtruncatemacro{\r}{2*\j+1}
            \draw[edge] (b\l.south) -- (c\j.north west);
            \draw[edge] (b\r.south) -- (c\j.north east);
        }
        % the path of the query x_0 = 2
        \node[cell, hot] at (a4) {$2$};
        \node[cell, hot] at (a5) {$3$};
        \node[cell, hot] at (b2) {$2$};
        \node[cell, hot] at (b3) {$8$};
        \node[cell, hot] at (c1) {$14$};
        \draw[hedge] (a4.south) -- (b2.north west);
        \draw[hedge] (a5.south) -- (b2.north east);
        \draw[hedge] (b2.south) -- (c1.north west);
        \draw[hedge] (b3.south) -- (c1.north east);

        % layer names and challenges
        \node[lab, anchor=east] at (-0.6, 0) {$g_0$ on $\Omega_0$};
        \node[lab, anchor=east] at (-0.6, -1.7) {$g_1$ on $\Omega_1$};
        \node[lab, anchor=east] at (-0.6, -3.4) {$g_2$ on $\Omega_2$};
        \node[lab, text=oc-green-9, anchor=west] at (\dx*7+0.5, -0.85) {fold, $\beta_0 = 2$};
        \node[lab, text=oc-green-9, anchor=west] at (\dx*7+0.5, -2.55) {fold, $\beta_1 = 5$};
        \node[lab, text=oc-gray-7, anchor=north] at (\dx*3.5, -3.9) {constant: the prover sends $c = 14$};
    \end{tikzpicture}
    \caption{The run of \Cref{ex:fri-run} as a tree. Every pair $\{x, -x\}$ of a layer (neighbouring cells, with the point written above each value) is folded into the single cell of $x^2$ in the next layer, so the layers form a binary tree. A query starting at $x_0 = 2$ climbs down one branch (orange): the verifier reads $g_0(\pm 2)$, then $g_1(\pm 4)$, and compares the last fold with $c$.}
    \label{fig:fri-tree}
\end{figure}


\begin{example}\label{ex:fri-cheat}
    Now let the prover try to pass off the evaluations of $T^5$,
    \begin{equation*}
        g_0 = \overline{T^5} = (1, 15, 4, 9, 16, 2, 13, 8),
    \end{equation*}
    which is not a codeword of $\mathsf{RS}[\mathbb{F}_{17}, \Omega, 4]$.
    Suppose it folds honestly, with the same challenges $\beta_0 = 2$ and $\beta_1 = 5$.
    Since $T^5 = T \cdot (T^2)^2$, we get $g_E = 0$ and $g_O = \overline{T^2}$, so $g_1 = \overline{2T^2} = (2, 15, 2, 15)$: the excess degree survived the fold.
    Folding again gives $g_2 = \overline{2T} = (2, 15)$ on $\Omega_2 = \{1, -1\}$, which is not constant.
    Whatever constant $c$ the prover sends, it differs from $g_2$ at one of the two points, and each query catches this with probability $1/2$.
    The prover could instead send a $g_1$ that \emph{is} a codeword, but then $g_1$ is inconsistent with $g_0$ on many pairs, and the first round consistency check catches it.
\end{example}

\begin{exercise}
    Why must the verifier send $\beta_i$ only \emph{after} receiving $g_i$? Describe a cheating strategy if all challenges were known in advance.
\end{exercise}

%--------------------------------
\subsubsection{Efficiency and Soundness}\label{subsubsec:fri-soundness}
%--------------------------------

\paragraph{Efficiency.}
The prover computes each fold in time linear in the size of the current domain, so its total work is $O(n + n/2 + \dots) = O(n)$ field operations.
The oracles $g_1, \dots, g_{\ell-1}$ have total length $n/2 + n/4 + \dots < n$, so the proof length is less than the length of the input.
Every repetition of the query phase reads $2$ values in each of the $\ell$ rounds, so the query complexity is $2\ell s = O(s \log n)$, and the verifier's work is of the same order.

\paragraph{Soundness.}
In the commit phase, by the second claim of \Cref{lemma:two-way-folding}, every round keeps a far function far, except with probability $\varepsilon_\delta$ over $\beta_i$.
If the prover sends an oracle $g_{i+1}$ that is close to the code, then $g_{i+1}$ must differ from the true fold of $g_i$ on many points, and the round consistency check detects this at a random point.
Either way, a far input leaves a trace of inconsistencies of density about $\delta$ along the path of a random query.
The following theorem makes this intuition precise.

\begin{theorem}[FRI soundness, informal]\label{thm:fri-soundness}
    Let $g_0$ be $\delta$-far from $\mathsf{RS}[\mathbb{F}, \Omega, k]$.
    Then for every prover, the FRI verifier with $s$ repetitions accepts with probability at most
    \begin{equation*}
        \ell \cdot \varepsilon_\delta + (1 - \delta)^s,
    \end{equation*}
    for every $\delta$ below the Johnson bound $1 - \sqrt{\rho}$ (with $\varepsilon_\delta$ as in \Cref{lemma:proximity-gap-simplified})~\cite{FOCS:BCIKS20}.
\end{theorem}

The first term is the probability that some challenge $\beta_i$ is unlucky, and it is negligible for a large field.
The second term is the probability that all $s$ query paths miss the inconsistencies, and it is what determines the number of queries in practice.

\begin{remark}
    The original analysis in~\cite[Theorem~3.3]{bensasson18} gives a weaker bound: the verifier accepts with probability at most $3n/\#\mathbb{F} + (1 - \min\{\delta, \delta_0\})^s$, where $\delta_0 \approx (1 - 3\rho)/4$ lies below the unique decoding radius.
    The improvement to the Johnson bound came later, first with DEEP-FRI~\cite{BGKS20} % TODO: bib key
    and then with the proximity gap theorem of~\cite{FOCS:BCIKS20}.
    Already~\cite[Conjecture~1.5]{bensasson18} conjectured that the bound holds up to the capacity $1 - \rho$; this is the conjecture discussed in \Cref{subsubsec:list-decoding}.
\end{remark}

\begin{exercise}
    Let $\rho = 1/4$, and ignore the first term of the soundness error in \Cref{thm:fri-soundness}.
    Find the smallest $s$ for which the remaining term is at most $2^{-100}$, using (a) the original bound from the remark above, (b) \Cref{thm:fri-soundness} with $\delta$ approaching the Johnson bound, and (c) the conjectured bound with $\delta$ approaching the capacity $1 - \rho$.

    \commentgray{Answer: $s \geq 100 / \log_2(1/(1-\delta))$, which gives about $1074$, $100$ and $50$ repetitions.}
\end{exercise}

\begin{exercise}
    Every query reads the pair $g_i(x_i), g_i(-x_i)$ in every round, and each value comes with a Merkle authentication path.
    How should the prover arrange the leaves of the Merkle tree of $g_i$ to halve the number of paths?
\end{exercise}

%--------------------------------
\subsubsection{FRI in the Big Picture}\label{subsubsec:fri-big-picture}
%--------------------------------

Let us now return to \Cref{fig:poly-iop-third-attempt}.
There, the input of the RS-IOPP is the virtual function $h = \sum_j \alpha^j \mathsf{q}_j$, which the prover never commits to.
FRI handles this without any change: whenever the verifier needs $h(\pm x_0)$ in the first round, it queries every committed oracle of \Cref{fig:poly-iop-third-attempt} at $\pm x_0$ and computes $h(\pm x_0)$ itself.
So the first layer of FRI comes for free: only the folds $g_1, \dots, g_{\ell-1}$ and the constant $c$ are new messages of the prover.

What remains is to turn the IOP into a SNARK.
This is done by the transformation of Ben-Sasson, Chiesa and Spooner~\cite{BCS16}: % TODO: bib key
\begin{algoen}
    \item every oracle is replaced by the root of a Merkle tree (\Cref{section:commitment-schemes}), and every query is answered with an authentication path;
    \item every challenge, including the query points $x_0$, is derived by hashing the transcript so far (\Cref{subsection:fiat-shamir}).
\end{algoen}
Both steps need only a hash function, so the resulting proof system needs no trusted setup and is plausibly post-quantum.
\Cref{fig:fri-pipeline} summarizes the whole path, which is the skeleton of every STARK.

\begin{figure}[ht]
    \centering
    \begin{tikzpicture}[
        box/.style={draw=oc-gray-7, thick, rounded corners=3pt, fill=oc-gray-0, align=center, minimum height=1.2cm, text width=2.0cm, font=\small},
        pipe/.style={-{Stealth[length=2.5mm]}, very thick, oc-blue-7},
        lab/.style={font=\scriptsize, align=center, text=oc-blue-9},
    ]
        \node[box] (a) at (0, 0) {Poly-IOP\\(e.g., PlonK, AIR)};
        \node[box] (b) at (4.4, 0) {IOP\\with one FRI run};
        \node[box, fill=oc-blue-1] (c) at (8.8, 0) {SNARK\\(STARK)};
        \draw[pipe] (a) -- node[lab, above] {bind, quotient,\\batch} node[lab, below] {\Cref{fig:poly-iop-third-attempt}} (b);
        \draw[pipe] (b) -- node[lab, above] {Merkle trees,\\Fiat-Shamir} node[lab, below] {BCS} (c);
    \end{tikzpicture}
    \caption{From a Poly-IOP to a SNARK using only a hash function. Polynomials are replaced by Reed-Solomon codewords and evaluation claims by quotients, which are batched and tested by a single run of FRI; the resulting IOP is compiled into a SNARK by committing to every oracle with a Merkle tree and deriving all challenges with Fiat-Shamir.}
    \label{fig:fri-pipeline}
\end{figure}

%--------------------------------
\subsubsection{Variants and Newer Protocols}\label{subsubsec:fri-variants}
%--------------------------------

Of course the protocol of \Cref{def:fri} is the simplest approach that has seen numerous improvements since the first introduction.
The most common change is the \emph{folding factor}.
Instead of pairs $\{x, -x\}$, one can fold whole cosets of size $2^\eta$ at once, using the map $x \mapsto x^{2^\eta}$ and reducing the degree by a factor of $2^\eta$ in one round.
The fold is then the value at $\beta_i$ of the polynomial of degree less than $2^\eta$ that interpolates $g_i$ on the coset, exactly as in \Cref{ex:fri-line}, which was the case $\eta = 1$.
\Cref{fig:fri-folding-factor} shows the benefit: a query reads slightly more values, but they all sit in one Merkle leaf, so far fewer authentication paths are sent.
The original paper uses $\eta = 2$~\cite{bensasson18}, and implementations typically use $\eta \in \{2, 3, 4\}$.
Another common change is to stop folding early, before reaching a constant, and to send the remaining polynomial of small degree in full, as the original paper does.

\begin{figure}[ht]
    \centering
    \begin{tikzpicture}[
        dot/.style={circle, fill=oc-blue-6, inner sep=1.5pt},
        hot/.style={circle, fill=oc-orange-7, inner sep=1.7pt},
        edge/.style={oc-gray-5, thick},
        lab/.style={font=\footnotesize, align=center},
    ]
        % 2-to-1 folding, three rounds
        \foreach \j in {0,...,7} { \node[dot] (l0\j) at (0.5*\j, 0) {}; }
        \foreach \j in {0,...,3} { \node[dot] (l1\j) at (1.0*\j+0.25, -0.8) {}; }
        \foreach \j in {0,1} { \node[dot] (l2\j) at (2.0*\j+0.75, -1.6) {}; }
        \node[hot] (l30) at (1.75, -2.4) {};
        \foreach \j in {0,...,3} { \pgfmathtruncatemacro{\l}{2*\j}\pgfmathtruncatemacro{\r}{2*\j+1}
            \draw[edge] (l0\l) -- (l1\j); \draw[edge] (l0\r) -- (l1\j); }
        \foreach \j in {0,1} { \pgfmathtruncatemacro{\l}{2*\j}\pgfmathtruncatemacro{\r}{2*\j+1}
            \draw[edge] (l1\l) -- (l2\j); \draw[edge] (l1\r) -- (l2\j); }
        \draw[edge] (l20) -- (l30); \draw[edge] (l21) -- (l30);
        % one query path
        \foreach \n in {l04, l05, l12, l13, l20, l21} { \node[hot] at (\n) {}; }
        \draw[oc-orange-7, very thick] (l04) -- (l12) -- (l21) -- (l30) (l05) -- (l12) (l13) -- (l21) (l20) -- (l30);
        \node[hot] at (l30) {};
        \node[lab] at (1.75, -3.2) {$2$-to-$1$ folding, $3$ rounds:\\$6$ values in $3$ Merkle openings};

        % 8-to-1 folding, one round
        \foreach \j in {0,...,7} { \draw[oc-orange-7, very thick] (6.0+0.5*\j, 0) -- (7.75, -2.4); \node[hot] (r\j) at (6.0+0.5*\j, 0) {}; }
        \node[hot] at (7.75, -2.4) {};
        \node[lab] at (7.75, -3.2) {$8$-to-$1$ folding, $1$ round:\\$8$ values in $1$ Merkle opening};
    \end{tikzpicture}
    \caption{Two ways to reduce the degree by a factor of $8$, and the values read by one query (orange). Folding pairs three times reads a pair in each of three layers, so it needs three Merkle openings. Folding a coset of size $8$ at once reads the whole coset in the first layer, which can be stored in a single leaf of the Merkle tree.}
    \label{fig:fri-folding-factor}
\end{figure}

The domains may also differ.
The original FRI is stated for binary fields, where squaring is not two-to-one; there, the domains are additive cosets, and the map $x \mapsto x^2$ is replaced by a polynomial that vanishes on a small subspace~\cite{bensasson18}.
Some approaches take even a step further: for example, Circle STARKs~\cite{EPRINT:HLP24} % TODO: bib key
work over $p = 2^{31} - 1$, where $p - 1$ has no large power of two as a factor, and use the points of the circle $x^2 + y^2 = 1$ instead, a group of the smooth order $p + 1 = 2^{31}$.

Finally, FRI is no longer the only RS-IOPP in use.
STIR~\cite{stir} % TODO: bib key for ePrint 2024/390
lowers the rate in every round in addition to the degree, so that later rounds need fewer queries, and it achieves this by performing the out-of-domain sampling and quotienting of \Cref{fig:poly-iop-second-attempt} inside every round.
WHIR~\cite{EPRINT:ACFY24b} % TODO: bib key
and Basefold~\cite{C:ZCF24} % TODO: bib key
carry the same ideas over to multilinear polynomials, which combines our considered approach with sum-check-based Poly-IOPs of \Cref{section:sumcheck}.

\subsection{Arithmetic Intermediate Representation (AIR)}\label{subsec:air}
Similar to SNARK, the STARK proof system involves intermediate stages, the first two of which are the \emph{arithmetic constraint system} and the \emph{IOP}.

\begin{definition}\label{def:air}
  The \emph{arithmetic intermediate representation (AIR)} or arithmetic internal representation, is a way of describing a computation in terms of an execution trace that satisfies a number of constraints induced by the correct evolution of the state.
\end{definition}

\begin{definition}\label{def:aet}
  Let $\mathbb{F}_q$ be the field of definition, and the computation describes the evolution of a state of $w$ registers for $\tau$ cycles. Then we define the \emph{algebraic execution trace (AET)} is the table of $\tau \times w$ field elements where every row describes the state of the system at the given point in time, and every column tracks the value of the given register.
\end{definition}

\begin{definition}\label{def:st-function}
  A \emph{state transition function} $\phi:\mathbb{F}_q^{w}\rightarrow\mathbb{F}_q^{w}$ determines the state at the next cycle as a function of the state at the previous cycle. Furthermore, a set of boundary conditions 
  \begin{equation*}
    \mathcal{B}\subseteq \mathbb{Z}_\tau \times\mathbb{Z}_w\times\mathbb{F}_q
  \end{equation*}
  fixes the values of some registers at some cycles: a triple $(i,j,e) \in \mathcal{B}$ means that register $j$ holds the value $e$ at cycle $i$.
\end{definition}

The \emph{computational integrity claim} consists of the state transition function and the boundary conditions. The \emph{witness} to this claim is the algebraic execution trace. The claim is \emph{true} if there is a witness $W\in \mathbb{F}_q^{\tau \times w}$ such that:
\begin{itemize}
  \item For every cycle, the state evolves correctly: $\forall i\in [\tau-1], \phi(W_{[i, :]})=W_{[i+1, :]}$
  \item All boundary conditions are satisfied: $\forall (i, j, e)\in \mathcal{B}, W_{[i, j]}=e$
\end{itemize}
Actually, the state transition function hides a lot of complexity and requires to be describable as low degree polynomials that are independent of the cycle.
However, this list of polynomials does not need to compute the next state from the current one; it merely needs to distinguish correct evolutions from incorrect ones. Specifically, the state transition function $\phi$ is represented by a list of polynomials $\mathbf{p}(x_0,\ldots,x_{w-1}, y_0,\ldots,y_{w-1})$ such that $\phi(\mathbf{x})=\mathbf{y}$ if and only if $\mathbf{p}(\mathbf{x},\mathbf{y})=0$. Say there are $m$ such state transition verification polynomials. Then the transition constraint become:
\begin{equation*}
  \forall i\in[\tau-1],\forall  \ell\in[m]: p_\ell(W_{[i, 0]},\ldots,W_{[i, w-1]},W_{[i+1, 0]},\ldots,W_{[i+1, w-1]})=0
\end{equation*}
This representation admits non-determinism: the verification polynomials may have much lower degree than any polynomial that computes the next state. For instance, the transition $y = x^{-1}$ requires computing $x^{q-2}$, but it is verified by the degree-$2$ polynomial $xy - 1$.

\subsubsection{Trace interpolation}
We saw that the arithmetic constraint system described above represents the computational integrity claim as a bunch of polynomials, where each such polynomial corresponds to a constraint. However, we need to extend terms of polynomials to the witness for transformation this constraint system into an IOP, and extending the notion of valid witnesses to witness polynomials.

Let $H = \langle h \rangle$ be a multiplicative subgroup of $\mathbb{F}_q^{\times}$ of order $\tau$, which we call the \emph{trace evaluation domain}. For simplicity, we assume that $\tau$ is a power of two; otherwise the trace is padded. The trace domain is smaller than the FRI evaluation domain $D$ by a factor exactly equal to the blowup factor $1/\rho$.

For every register $j \in [w]$, the \emph{trace polynomial} $t_j \in \mathbb{F}_q[T]$ is the unique polynomial of degree less than $\tau$ such that
\begin{equation*}
  t_j(h^i) = W_{[i,j]} \quad \forall i \in [\tau].
\end{equation*}
In practice, it is computed with the inverse NTT (\Cref{section:ntt}). The trace polynomials
$t_0, \ldots, t_{w-1}$ represent the algebraic execution trace in terms of univariate polynomials.

Translating the computational integrity claim to the trace polynomials, we get:
\begin{itemize}
  \item All boundary constraints are satisfied: $\forall (i,j,e)\in \mathcal{B}, t_j(h^i)=e$.
  \item For all cycles, all transition constraints are satisfied: $\forall i\in [\tau-1],\forall \ell\in[m]$
  \begin{equation*}
    p_\ell(t_0(h^i),\dots,t_{w-1}(h^i),t_0(h^{i+1}),\dots,t_{w-1}(h^{i+1}))=0
  \end{equation*}
\end{itemize}
If we look at the left-hand side of the equation, we can see that it corresponds to a univariate polynomial $p_\ell(t_0(T),\dots,t_{w-1}(T),t_0(h\cdot T),\dots,t_{w-1}(h\cdot T))$. Which just means that all $m$ of these transition polynomials evaluate to $0$ in $H$.

Then we can build the following high-level protocol:
This observation gives rise to the following high-level protocol:
\begin{algoen}
  \item The prover commits to the trace polynomials $t_0(T), \ldots, t_{w-1}(T)$.
  \item The verifier checks that $t_j(h^i) = e$ for all $(i,j,e) \in \mathcal{B}$.
  \item The prover commits to the transition polynomials
  \begin{equation*}
    c_\ell(T) := p_\ell\big(t_0(T), \ldots, t_{w-1}(T), t_0(h \cdot T), \ldots, t_{w-1}(h \cdot T)\big),
    \quad \ell \in [m].
  \end{equation*}
  \item The verifier checks that $c_\ell(T)$ and $t_j(T)$ are correctly related by:
  \begin{itemize}
    \item choosing a random point $z \leftarrowS \mathbb{F}_q^{\times}$;
    \item querying the values $t_j(z)$ and $t_j(h \cdot z)$ $\forall j \in [w]$;
    \item evaluating the transition verification polynomials $p_\ell$ at
    \begin{equation*}
      \big(t_0(z), \ldots, t_{w-1}(z), t_0(h \cdot z), \ldots, t_{w-1}(h \cdot z)\big)
    \end{equation*}
    \item querying the values $c_\ell(z)$;
    \item checking that the values obtained in the previous two steps match.
  \end{itemize}
  \item The verifier checks that the transition polynomials $c_\ell(T)$ evaluate to zero on
  $\{h^i : i \in [\tau-1]\}$.
\end{algoen}

In fact, the commitment to the transition polynomials can be omitted. Instead, the verifier uses
the values of $t_j$ at $z$ and $h \cdot z$ to compute the value of $c_\ell(T)$ at the one point
needed to verify that $c_\ell(T)$ evaluates to $0$ on $\{h^i : i \in [\tau-1]\}$.

We could use the protocol above if it did not have a flaw: the verifier cannot perform steps 2 and 5, because he sees the values of the polynomials only on the domain $D$, while these checks concern the trace domain $H$. That is why we use FRI to simulate these evaluation checks, with a trick similar to the one used for KZG polynomial commitments. By \Cref{th:root-div}, $g(x) = y$ if and only if $(T - x) \mid g(T) - y$. Hence an evaluation check can be simulated by:
\begin{itemize}
  \item subtracting the $y$-coordinate;
  \item dividing out the \emph{zerofier}, which is the minimal polynomial that vanishes at the $x$-coordinate;
  \item proving with FRI that the resulting quotient has a bounded degree.
\end{itemize}

This procedure happens twice for the STARK polynomials: (1) first, applied to the trace polynomials to show satisfaction of the boundary constraints and (2) second, applied to the transition polynomials to show that the transition constraints are satisfied. We call the resulting lists of quotient polynomials the \emph{boundary quotients} and the \emph{transition quotients} respectively.

Both kinds of constraints are instances of a single, more general notion.

\begin{definition}[Generalized AIR constraint]\label{def:gen-air}
  A \emph{generalized AIR constraint} is given by two polynomials:
  \begin{itemize}
    \item The \emph{numerator} determines which equation between elements of the algebraic execution
    trace holds, in a manner that is independent of the cycle. For transition constraints, the numerator
    is exactly the transition polynomial $c_\ell(T)$. For a boundary constraint $(i,j,e) \in \mathcal{B}$,
    the numerator is simply $t_j(T) - e$.
    \item The \emph{denominator} is a zerofier that determines where the equation is supposed to hold,
    by vanishing in those points. For transition constraints, this zerofier vanishes on all points of the
    trace evaluation domain except the last. For boundary constraints, it vanishes only on the boundary points.
  \end{itemize}
  The constraint is satisfied if and only if the numerator is divisible by the denominator, that is,
  if the \emph{quotient} of the numerator by the denominator is a polynomial.
\end{definition}

For every register $j \in [w]$, let $\mathcal{B}_j := \{h^i : (i,j,e) \in \mathcal{B}\}$ be its set of
boundary points, let the \emph{boundary interpolant} $I_j$ be the polynomial of degree less than
$|\mathcal{B}_j|$ such that $I_j(h^i) = e$ for all $(i,j,e) \in \mathcal{B}$, and let
$Z_{\mathcal{B}_j}(T) := \prod_{x \in \mathcal{B}_j} (T - x)$ be the \emph{boundary zerofier}.
The \emph{transition zerofier} is
\begin{equation*}
  Z(T) := \prod_{i \in [\tau-1]} (T - h^i) = \frac{T^\tau - 1}{T - h^{\tau-1}},
\end{equation*}
where the second equality follows from \Cref{lemma:vanishing-polynomial}. The quotients are
\begin{equation*}
  q^{\mathcal{B}}_j(T) := \frac{t_j(T) - I_j(T)}{Z_{\mathcal{B}_j}(T)}, \quad j \in [w],
  \qquad
  q^{\mathcal{T}}_\ell(T) := \frac{c_\ell(T)}{Z(T)}, \quad \ell \in [m].
\end{equation*}

The degree bounds are $\deg q^{\mathcal{B}}_j < \tau - |\mathcal{B}_j|$ and, if $p_\ell$ has degree $d_\ell$, $\deg q^{\mathcal{T}}_\ell \leq (d_\ell - 1)(\tau - 1)$. As a result, we can summarize the discussion as follows.

\begin{proposition}\label{prop:stark-quotients}
  The trace $W$ satisfies all boundary and transition constraints if and only if all boundary quotients $q^{\mathcal{B}}_j$ and all transition quotients $q^{\mathcal{T}}_\ell$ are polynomials.
\end{proposition}

\begin{proof}
  A quotient is a polynomial if and only if its numerator vanishes at every root of its denominator (\Cref{th:root-div}, applied to each root in turn). For the boundary quotients this means $t_j(h^i) = e$ for all $(i,j,e) \in \mathcal{B}$, and for the transition quotients it means $c_\ell(h^i) = 0$ for all $i \in [\tau-1]$, which are exactly the constraints.
\end{proof}

The trace polynomial $t_j$ appears in both types of quotients, while $I_j$ and $Z_{\mathcal{B}_j}$ are public. Therefore, the trace can be reconstructed from the boundary quotient,
\begin{equation*}
  t_j(T) = q^{\mathcal{B}}_j(T) \cdot Z_{\mathcal{B}_j}(T) + I_j(T),
\end{equation*}
and the prover commits to the boundary quotients instead of the trace polynomials. After opening $q^{\mathcal{B}}_j$ at the points $x$ and $h \cdot x$, the verifier computes $t_j$ itself, then derives $c_\ell$ from it, and finally derives $q^{\mathcal{T}}_\ell$ from $c_\ell$.

\begin{remark}
  The zerofiers $Z$ and $Z_{\mathcal{B}_j}$ vanish on points of $H$, so the quotients cannot be evaluated there. Since $H = \langle h \rangle$ is a subgroup of $\langle \omega \rangle$, we take the FRI domain to be a coset $D = c \cdot \langle \omega \rangle$ with $c \notin \langle \omega \rangle$, which does not intersect $H$. Moreover, $h = \omega^{1/\rho}$, so $h \cdot x \in D$ for every $x \in D$. FRI works on a coset in exactly the same way: $-x \in D$ for every $x \in D$, and since $(cx)^2 = c^2 x^2$, squaring maps $D$ onto the coset $c^2 \cdot \langle \omega^2 \rangle$ of half the size.
\end{remark}

\subsubsection{Nonlinear Combination}

Instead of running FRI separately for every quotient, the prover combines all of them into a single polynomial. For simplicity, we assume from now on that all transition verification polynomials have degree at most $2$, as in~\cite{szepieniec2021anatomy}. Then every quotient has degree less than $\tau$, and we run FRI with the degree bound $k := \tau$, so that $n = \tau/\rho$.

Let $\mathcal{Q}$ denote the list of all boundary and transition quotients, and for $q \in \mathcal{Q}$ let $d_q$ be its degree bound: $d_q = \tau - |\mathcal{B}_j| - 1$ for $q = q^{\mathcal{B}}_j$, and $d_q = (d_\ell - 1)(\tau - 1)$ for $q = q^{\mathcal{T}}_\ell$. After receiving random weights $\alpha_q, \beta_q \leftarrowS \mathbb{F}_q$ from the verifier, the prover computes
\begin{equation*}
  g(T) := \sum_{q \in \mathcal{Q}} \big( \alpha_q \cdot q(T) + \beta_q \cdot T^{\,k-1-d_q} \cdot q(T) \big).
\end{equation*}

Each quotient enters the sum twice. The first term alone would not be enough: a quotient whose degree exceeds its own bound $d_q$ but is still less than $k$ would not affect the degree of $g$ and would go unnoticed. The second term shifts $q$ so that its degree bound becomes exactly $k-1$: if $\deg q > d_q$, this term has degree at least $k$. Hence, with high probability over the choice of the weights, $\deg g < k$ if and only if $\deg q \leq d_q$ for every $q \in \mathcal{Q}$. Likewise, if some quotient is not a polynomial at all, then with high probability the codeword of $g$ on $D$ is far from $\mathsf{RS}_q[D,k]$. Following~\cite{szepieniec2021anatomy}, we call $g$ a \emph{nonlinear combination}, because of the factors $T^{\,k-1-d_q}$. It therefore suffices to run FRI once, on $g$.

\subsection{The STARK Protocol}\label{subsec:stark-protocol}

We can now put everything together. Note that the random point $z$ from the high-level protocol is no
longer needed: its role is played by the points that the verifier queries in the query phase of FRI.

\begin{definition}[STARK protocol]\label{def:stark}
  The AIR $\big(\tau, w, \{p_\ell\}_{\ell \in [m]}, \mathcal{B}\big)$, the domains $H$ and $D$, and the
  degree bound $k$ are public; the prover knows a trace $W$.
  \begin{algoen}
    \item The prover computes the trace polynomials $t_j$ and the boundary quotients $q^{\mathcal{B}}_j$,
    and sends the Merkle roots of the codewords of $q^{\mathcal{B}}_j$ on $D$ for all $j \in [w]$.
    \item The verifier sends random weights $\alpha_q, \beta_q \leftarrowS \mathbb{F}_q$ for all
    $q \in \mathcal{Q}$.
    \item The prover computes the transition quotients $q^{\mathcal{T}}_\ell$ and the nonlinear
    combination $g$, and runs the commit phase of FRI (\Cref{def:fri}) with $f_0 := g$.
    \item The parties run the query phase of FRI. For every point $x \in D$ at which the verifier reads
    $f_0$, the prover additionally opens $q^{\mathcal{B}}_j(x)$ and $q^{\mathcal{B}}_j(h \cdot x)$ for all
    $j \in [w]$, together with their authentication paths.
    \item For every such $x$, the verifier checks the authentication paths, recovers $t_j(x)$ and
    $t_j(h \cdot x)$ from the boundary quotients, computes $c_\ell(x)$,
    $q^{\mathcal{T}}_\ell(x) = c_\ell(x)/Z(x)$ and $g(x)$, and checks that $g(x) = f_0(x)$.
  \end{algoen}
  The verifier accepts if and only if all these checks and all checks of FRI pass.
\end{definition}

\begin{figure}[h!]
    \centering
    \scalebox{0.74}{
    \begin{tikzpicture}[
        arr/.style   = {-{Stealth[length=2.5mm]}, line width=0.4mm},
        act/.style   = {draw=oc-blue-7, fill=oc-blue-1, rounded corners=2pt,
                        line width=0.4mm, inner sep=4pt, font=\small},
        data/.style  = {draw=gray, rounded corners=2pt, line width=0.4mm,
                        inner sep=4pt, font=\small},
        lbl/.style   = {font=\scriptsize, inner sep=2pt},
        phase/.style = {draw=gray, dashed, rounded corners=8pt, line width=0.4mm},
        proc/.style  = {draw=gray, fill=gray!20!white, rounded corners=2pt,
                        line width=0.4mm, inner sep=4pt, font=\small},
        check/.style = {draw=black, rounded corners=2pt, line width=0.6mm,
                        inner sep=4pt, font=\small},
    ]
        % ---------------- Prover row
        \draw[phase] (-0.8,1.05) rectangle (13.1,-2.3);
        \node[font=\scriptsize\scshape, text=gray, anchor=north east] at (13.0,0.95) {Prover};
        \node[inner sep=0pt, align=center] at (-1.6,-0.6)
            {\includegraphics[width=1.0cm]{contents/images/common/prover.png}\\ $\mathcal{P}$};

        \node[data] (W)  at (-0.15,0)   {$W$};
        \node[data] (tj) at (2.3,0)     {$t_j$};
        \node[act]  (qB) at (6.7,0)     {$q^{\mathcal{B}}_j$};
        \node[data] (cl) at (4.6,-1.55) {$c_\ell$};
        \node[data] (qT) at (6.7,-1.55) {$q^{\mathcal{T}}_\ell$};
        \node[act]  (g)  at (9.2,-0.78) {$g$};
        \node[proc] (fri) at (11.8,-0.78) {FRI (\Cref{fig:fri-scheme})};

        \draw[arr] (W) -- node[lbl, above] {inverse NTT} (tj);
        \draw[arr] (tj) -- node[lbl, above] {$(t_j - I_j)/Z_{\mathcal{B}_j}$} (qB);
        \draw[arr, rounded corners=4pt] (tj.south) |- node[lbl, pos=0.75, below] {$p_\ell$} (cl.west);
        \draw[arr] (cl) -- node[lbl, above] {$\div Z$} (qT);
        \draw[arr] (qB.east) -- (g.north west);
        \draw[arr] (qT.east) -- (g.south west);
        \draw[arr] (g) -- (fri);
        \node[lbl, above=1mm of g] {$\alpha_q, \beta_q$};

        % ---------------- Verifier row
        \draw[phase] (-0.8,-2.75) rectangle (13.1,-5.35);
        \node[font=\scriptsize\scshape, text=gray, anchor=north east] at (13.0,-2.85)
            {Verifier, at every point $x$ queried by FRI};
        \node[inner sep=0pt, align=center] at (-1.6,-4.05)
            {\includegraphics[width=1.0cm]{contents/images/common/verifier.png}\\ $\mathcal{V}$};

        \node[data] (v1) at (1.05,-3.9) {$q^{\mathcal{B}}_j(x),\ q^{\mathcal{B}}_j(h x)$};
        \node[data] (v2) at (4.6,-3.9)  {$t_j(x),\ t_j(h x)$};
        \node[data] (v3) at (6.9,-3.9)  {$c_\ell(x)$};
        \node[data] (v4) at (8.8,-3.9)  {$q^{\mathcal{T}}_\ell(x)$};
        \node[check] (v5) at (11.5,-3.9) {$g(x) \stackrel{?}{=} f_0(x)$};

        \draw[arr] (v1) -- (v2);
        \draw[arr] (v2) -- (v3);
        \draw[arr] (v3) -- (v4);
        \draw[arr] (v4) -- (v5);
        \draw[arr, rounded corners=4pt] (v1.south) -- ++(0,-0.55) -| (v5.south);
    \end{tikzpicture}
    }
    \caption{Overview of the STARK protocol. The prover commits only to the highlighted polynomials:
    the boundary quotients $q^{\mathcal{B}}_j$ and, through FRI, the nonlinear combination $g$. The trace
    polynomials $t_j$, the transition polynomials $c_\ell$ and the transition quotients
    $q^{\mathcal{T}}_\ell$ are never committed. At every point queried by FRI, the verifier recomputes
    them from the opened values of the boundary quotients and compares the result with the first layer
    of FRI. Adapted from~\cite{szepieniec2021anatomy}.}
    \label{fig:stark-scheme}
\end{figure}

If the trace is correct, the honest prover is always accepted. If it is not, then by
\Cref{prop:stark-quotients} at least one quotient is not a polynomial of bounded degree, so with high
probability the codeword of $g$ is far from $\mathsf{RS}_q[D,k]$, and FRI rejects. The prover cannot
avoid this by running FRI on a different function that is close to a codeword: step 5 ties the values
of $f_0$ at the queried points to the committed boundary quotients.

\begin{remark}
  As with FRI, the whole protocol is made non-interactive using the Fiat--Shamir transformation
  (\Cref{subsection:fiat-shamir}). Practical STARKs also add several details that we omit: randomization
  of the trace for zero-knowledge, and additional checks that improve soundness, such as DEEP.
\end{remark}

\end{document}